\documentclass[../DoAn.tex]{subfiles}
\begin{document}

Trong thời đại công nghệ phát triển, chuyển đổi số trong giáo dục đại học trở thành một vấn đề vô cùng quan trọng đối với việc nâng cao hiệu quả giảng dạy, trải nghiệm và chất lượng dịch vụ cho toàn bộ nhà trường nói chung và cá nhân từng sinh viên nói riêng. Trong đó, quy trình đăng ký, phân bổ phòng và quản lý ký túc xá là một bài toán thực tế với số lượng người dùng lớn, tồn đọng vô cùng nhiều quy định và thường xuyên thay đổi theo từng năm học, thậm chí là theo từng tháng. Cùng với đó, trải nghiệm trực quan của sinh viên vẫn còn là một hạn chế lớn do thiếu các công cụ hỗ trợ quan sát và khám phá không gian phòng ở một cách chân thực trước khi đăng ký. Việc chỉ cung cấp thông tin dưới dạng văn bản hoặc một số hình ảnh tĩnh chưa phản ánh đầy đủ điều kiện thực tế của phòng, gây khó khăn cho sinh viên trong quá trình lựa chọn và làm giảm tính minh bạch của dịch vụ ký túc xá.

\section{Đặt vấn đề}
\label{section:1.1}

Quản lý và phân bổ phòng ký túc xá là một trong những hoạt động quan trọng tại các trường đại học vào đầu mỗi kỳ tuyển sinh, ảnh hưởng trực tiếp đến điều kiện sinh hoạt, học tập và trải nghiệm của sinh viên trong suốt quá trình học tập. Với số lượng sinh viên ngày càng lớn cả trong và ngoài nước, việc tổ chức đăng ký ký túc xá không còn chỉ là bài toán quản lý dữ liệu mà trở thành một bài toán cần đảm bảo tính công bằng, minh bạch và hiệu quả trong khâu vận hành.

Hiện nay, đa số hệ thống ký túc xá của nhiều trường đại học trên toàn Việt Nam vẫn sử dụng những biện pháp thủ công như nhập liệu, lưu trữ hồ sơ giấy hay thống kê hoặc lập báo cáo bằng Excel với số lượng giấy tờ, hồ sơ khổng lồ. Điều đó dẫn đến sự khó khăn trong việc quản lý dữ liệu, theo dõi trạng thái hồ sơ, đảm bảo tính minh bạch của quá trình phân phòng cũng như kéo dài thời gian hoàn thành từng tác vụ bằng sức người. Bên cạnh đó, sinh viên cũng hạn chế trong việc trải nghiệm và hình dung được không gian mà mình sẽ gắn bó trong quãng thời gian đầu khi bước vào đại học bởi thông tin cung cấp chỉ qua những bài đăng trên website hoặc trang mạng xã hội của nhà trường.


Tính minh bạch của quá trình xét duyệt và phân bổ phòng cũng là một vấn đề vô cùng nan giải đối với toàn bộ hệ thống hiện nay. Sinh viên thường chỉ nhận được kết quả cuối cùng qua Gmail hoặc một kênh liên lạc cố định mà không có nhiều thông tin về cách thức đánh giá hồ sơ, các tiêu chí ưu tiên được áp dụng hay lý do dẫn đến kết quả phân bổ. Điều này có thể tạo ra những thắc mắc hoặc khiếu nại trong quá trình vận hành, đồng thời làm tăng khối lượng công việc giải đáp của cán bộ quản lý.


Bên cạnh các vấn đề về quản lý, trải nghiệm của sinh viên cũng đang đặt ra nhiều yêu cầu mới đối với các hệ thống ký túc xá hiện đại. Sinh viên ngày nay có xu hướng thực hiện hầu hết các hoạt động trên thiết bị di động và mong muốn có thể đăng ký, theo dõi trạng thái hồ sơ, nhận thông báo và tra cứu thông tin một cách nhanh chóng. Đặc biệt, trước khi quyết định đăng ký, sinh viên thường có nhu cầu tìm hiểu môi trường sống, cơ sở vật chất và điều kiện phòng ở để đưa ra lựa chọn phù hợp. Tuy nhiên, phần lớn các hệ thống hiện nay chỉ cung cấp thông tin dưới dạng văn bản hoặc hình ảnh đơn lẻ, chưa cho phép người dùng có cái nhìn trực quan về không gian thực tế của phòng ở.

Việc thiếu các công cụ hỗ trợ trực quan khiến sinh viên khó hình dung điều kiện sinh hoạt trước khi đăng ký, đồng thời làm giảm mức độ tương tác và trải nghiệm của người dùng đối với hệ thống. Trong bối cảnh công nghệ thực tế ảo, ảnh 360° và mô hình không gian số ngày càng trở nên phổ biến, việc ứng dụng các công nghệ này vào lĩnh vực quản lý ký túc xá có thể mang lại nhiều giá trị thực tiễn. Sinh viên có thể tham quan phòng ở từ xa, quan sát bố trí không gian, cơ sở vật chất và môi trường xung quanh trước khi đưa ra quyết định đăng ký, từ đó nâng cao tính chủ động và mức độ hài lòng đối với dịch vụ ký túc xá.

Từ những vấn đề trên, có thể thấy nhu cầu xây dựng một hệ thống quản lý ký túc xá hiện đại là cần thiết. Hệ thống không chỉ hỗ trợ số hóa quy trình đăng ký và phân bổ phòng, nâng cao tính minh bạch trong quản lý mà còn hướng tới cải thiện trải nghiệm người dùng thông qua các công cụ trực quan hóa phòng ở và khả năng tương tác trên nền tảng số. 


\section{Mục tiêu và phạm vi đề tài}
\label{section:1.2}

\subsection{Mục tiêu nghiên cứu}

Mục tiêu chung của đề tài là nghiên cứu, thiết kế và xây dựng hệ thống E-Dormitory nhằm hỗ trợ số hóa quy trình quản lý ký túc xá, từ khâu đăng ký lưu trú, quản lý thông tin sinh viên, quản lý phòng ở đến xét duyệt và phân bổ chỗ ở. Hệ thống hướng tới việc giảm khối lượng thao tác thủ công, nâng cao khả năng theo dõi và kiểm soát dữ liệu, đồng thời tăng tính minh bạch, công bằng và thuận tiện trong quá trình phục vụ sinh viên.

Để đạt mục tiêu trên, đề tài cần thực hiện một số nhóm việc có quan hệ chặt chẽ với nhau trong toàn bộ vòng đời vận hành của hệ thống. Trước hết, đề tài xây dựng tầng backend làm nền tảng xử lý nghiệp vụ, bao gồm các chức năng quản lý sinh viên, ký túc xá, phòng ở, hồ sơ đăng ký và các dữ liệu liên quan đến quá trình xét duyệt. Trên nền tảng dữ liệu đó, hệ thống tiếp tục xây dựng cơ chế phân bổ phòng dựa trên các tiêu chí ưu tiên, chỉ tiêu phân bổ theo từng nhóm đối tượng và sức chứa thực tế của từng phòng, nhằm hỗ trợ cán bộ quản lý ra quyết định một cách có căn cứ và dễ kiểm tra.

Song song với phần xử lý nghiệp vụ, đề tài xây dựng giao diện web cho hai nhóm người dùng chính là sinh viên và cán bộ quản lý. Đối với sinh viên, giao diện web hỗ trợ tra cứu thông tin ký túc xá, gửi đăng ký, theo dõi trạng thái hồ sơ và xem kết quả phân bổ. Đối với cán bộ quản lý, giao diện web hỗ trợ quản trị dữ liệu, xét duyệt hồ sơ, theo dõi tình trạng phòng, thực hiện phân bổ và khai thác các báo cáo vận hành cần thiết. Bên cạnh nền tảng web, đề tài cũng phát triển ứng dụng di động dành cho sinh viên để đáp ứng thói quen truy cập thường xuyên trên thiết bị cá nhân, hỗ trợ đăng nhập, tra cứu, nhận thông báo và cập nhật trạng thái một cách thuận tiện hơn.

Ngoài các chức năng quản lý và đăng ký, đề tài hướng tới cải thiện trải nghiệm lựa chọn chỗ ở thông qua các công cụ trực quan hóa thông tin phòng và ký túc xá. Sinh viên có thể xem ảnh phòng, tham quan không gian qua ảnh hoặc tour 360°, đồng thời tiếp cận các mô hình trực quan giúp hình dung rõ hơn về bố trí, cơ sở vật chất và môi trường sinh hoạt trước khi đăng ký. Hệ thống cũng cần xây dựng cơ chế thông báo, theo dõi trạng thái, tổng hợp báo cáo và tổ chức dữ liệu vận hành để hỗ trợ cả sinh viên lẫn cán bộ quản lý trong quá trình sử dụng thực tế.

Với các mục tiêu trên, đề tài không chỉ dừng lại ở việc mô tả một mô hình quản lý mang tính học thuật mà được phát triển theo hướng full-stack, bao gồm backend, giao diện web, ứng dụng di động, cơ sở dữ liệu, cơ chế thời gian thực và các thành phần hỗ trợ triển khai. Định hướng này giúp hệ thống có khả năng chạy thử trong môi trường thực tế, qua đó đánh giá được tính khả thi của giải pháp cả ở góc độ kỹ thuật lẫn nghiệp vụ.


\subsection{Phạm vi đề tài}
Đề tài tập trung nghiên cứu, thiết kế và triển khai hệ thống E-Dormitory -- hệ thống quản lý và tiện ích ký túc xá cho sinh viên tích hợp phân bổ phòng tự động và trực quan hóa không gian. Phạm vi của đề tài bao gồm việc xây dựng các thành phần phần mềm chính của một hệ thống full-stack có thể vận hành thử nghiệm, trong đó backend đảm nhiệm xử lý nghiệp vụ và cung cấp API, giao diện web phục vụ sinh viên và cán bộ quản lý, ứng dụng di động phục vụ sinh viên, còn cơ sở dữ liệu lưu trữ các thông tin cốt lõi về người dùng, ký túc xá, phòng ở, đăng ký, phân bổ và thông báo.

Về nghiệp vụ, hệ thống phục vụ hai nhóm người dùng chính là sinh viên và cán bộ quản lý. Sinh viên có thể tra cứu thông tin ký túc xá, xem thông tin phòng, gửi đăng ký lưu trú, theo dõi trạng thái hồ sơ, nhận thông báo và xem kết quả phân bổ. Cán bộ quản lý có thể quản lý dữ liệu sinh viên, ký túc xá và phòng ở, tiếp nhận và xét duyệt hồ sơ đăng ký, cấu hình các tiêu chí phục vụ phân bổ, thực hiện phân phòng, theo dõi tình trạng sử dụng phòng và khai thác các báo cáo thống kê cơ bản.

Về cơ chế phân bổ phòng, đề tài tập trung vào hướng tiếp cận dựa trên tiêu chí ưu tiên, quota và sức chứa thực tế của phòng. Hệ thống không thay thế hoàn toàn vai trò của cán bộ quản lý mà đóng vai trò hỗ trợ tự động hóa một phần quy trình, giúp tạo danh sách xét duyệt, đề xuất hoặc ghi nhận kết quả phân bổ và cung cấp dữ liệu để kiểm tra lại khi cần thiết. Cách tiếp cận này phù hợp với bối cảnh ký túc xá thường có nhiều quy định vận hành thay đổi theo từng năm học, từng nhóm sinh viên hoặc từng đợt đăng ký.

Một nội dung quan trọng khác trong phạm vi đề tài là xây dựng các chức năng trực quan hóa thông tin ký túc xá. Sinh viên có thể xem thông tin chi tiết về cơ sở vật chất, hình ảnh phòng, ảnh hoặc tour 360° và mô hình trực quan trước khi đăng ký, qua đó có thêm cơ sở để lựa chọn chỗ ở phù hợp. Các chức năng thông báo, theo dõi trạng thái hồ sơ và báo cáo vận hành cũng được đưa vào phạm vi triển khai ở mức đáp ứng nhu cầu quản lý cơ bản và hỗ trợ quá trình chạy thử hệ thống.

Ngoài phạm vi đề tài là các chức năng tích hợp sâu với hệ thống đào tạo, quản lý điểm hoặc tuyển sinh, cũng như các bài toán phân tích dữ liệu nâng cao, dự báo nhu cầu bằng trí tuệ nhân tạo, thanh toán trực tuyến quy mô đầy đủ hoặc triển khai hạ tầng lớn cho toàn bộ trường đại học. Đây là những hướng phát triển tiềm năng trong tương lai khi hệ thống được đánh giá, hoàn thiện và đưa vào khai thác thực tế ở quy mô rộng hơn.
\section{Định hướng giải pháp}

\label{section:1.3}
Giải pháp được xây dựng dựa trên bốn định hướng chính: phân tách rõ ràng giữa các tầng chức năng, cho phép cấu hình chính sách linh hoạt mà không cần can thiệp mã nguồn, nâng cao trải nghiệm lựa chọn chỗ ở của sinh viên và hỗ trợ vận hành trên nhiều nền tảng.

Về kiến trúc, hệ thống được tổ chức theo mô hình ba tầng gồm giao diện, dịch vụ và dữ liệu, giúp tách biệt các chức năng xử lý nghiệp vụ, quản lý dữ liệu và tương tác người dùng, từ đó nâng cao khả năng bảo trì và mở rộng. Cơ chế mô phỏng và thử nghiệm chính sách trước khi áp dụng thực tế giúp cán bộ quản lý đánh giá hiệu quả các phương án vận hành một cách an toàn.

Hệ thống hỗ trợ sinh viên tìm hiểu và trải nghiệm phòng ở trước khi đăng ký thông qua hình ảnh thực tế, không gian 360° và thông tin chi tiết về cơ sở vật chất, giúp đưa ra quyết định phù hợp hơn so với hình thức đăng ký truyền thống.

Về công nghệ, hệ thống sử dụng Node.js và Express cho tầng dịch vụ, MongoDB Atlas cho lưu trữ dữ liệu, Redis hỗ trợ các chức năng thời gian thực, giao diện web phát triển bằng EJS kết hợp Bootstrap, đồng thời cung cấp ứng dụng di động xây dựng trên React Native và Expo nhằm đảm bảo trải nghiệm thống nhất trên web, Android và iOS. Hệ thống còn áp dụng cơ chế xác thực, phân quyền và kiểm thử nhiều lớp nhằm đảm bảo tính an toàn và ổn định trong vận hành.
\section{Bố cục đồ án}

\label{section:1.4}
Phần còn lại của báo cáo được tổ chức thành năm chương tiếp theo, hình thành một tiến trình nghiên cứu xuyên suốt từ khảo sát bài toán thực tiễn, phân tích yêu cầu, đề xuất giải pháp, thiết kế và triển khai hệ thống, đến trình bày các giải pháp, những đóng góp của đồ án và đánh giá kết quả đạt được.


Chương 2 — Khảo sát và phân tích yêu cầu nghiên cứu thực trạng quản lý và phân bổ phòng ký túc xá tại các cơ sở giáo dục, khảo sát các hệ thống đang được sử dụng trong và ngoài nước, từ đó xác định các hạn chế còn tồn tại, yêu cầu nghiệp vụ và khoảng trống mà đề tài hướng tới giải quyết.

Chương 3 — Công nghệ sử dụng và kiến trúc tổng thể hệ thống trình bày các công nghệ, nền tảng và công cụ được lựa chọn, đồng thời mô tả kiến trúc tổng thể của E-Dormitory và cách các thành phần phối hợp để đáp ứng yêu cầu nghiệp vụ.

Chương 4 — Thiết kế hệ thống đi sâu vào các quyết định tổ chức phần mềm ở mức thiết kế, giải thích lý do lựa chọn mô hình kiến trúc, cách các thành phần được phân tầng và phụ thuộc lẫn nhau, cũng như cấu trúc nội tại của các gói nghiệp vụ cốt lõi.

Chương 5 — Các đóng góp và giải pháp đề xuất trình bày những đóng góp chính của đồ án, trong đó mỗi đóng góp trả lời bốn câu hỏi cốt lõi: vấn đề cần giải quyết, đối tượng hưởng lợi, điểm khác biệt so với giải pháp hiện có và giá trị mang lại. Ba đóng góp được trình bày gồm: giải pháp trực quan hóa không gian phòng ký túc xá (Mục 5.1), giải pháp phân bổ phòng tự động dựa trên quota (Mục 5.2) và giải pháp kiến trúc phân bổ linh hoạt theo năm học (Mục 5.3).

Chương 6 — Kết luận và hướng phát triển tổng kết quá trình nghiên cứu, đánh giá kết quả đạt được, phân tích hạn chế còn tồn tại và đề xuất các định hướng hoàn thiện, mở rộng hệ thống nhằm đáp ứng yêu cầu vận hành thực tế ở quy mô toàn trường.

Ngoài các chương chính, báo cáo còn bao gồm phần Tài liệu tham khảo và Phụ lục cung cấp sơ đồ cơ sở dữ liệu, tài liệu API, hướng dẫn cài đặt và triển khai hệ thống.

Từ mục tiêu và phạm vi đã xác định trong chương này, nội dung tiếp theo của báo cáo chuyển sang khảo sát hiện trạng và phân tích yêu cầu. Chương 2 sẽ làm rõ bối cảnh quản lý ký túc xá, các khó khăn trong quy trình đăng ký và phân bổ phòng hiện nay, đồng thời xác định những yêu cầu chức năng và phi chức năng làm cơ sở cho việc thiết kế, triển khai hệ thống ở các chương sau.


\end{document}
